\section{Preamble}

Software systems are becoming increasingly autonomous. 
They are no longer limited to reporting information or 
executing narrowly specified commands. Instead, 
emerging software systems are expected to act as 
assistants: they customize behavior for individual users, 
carry out multi-step tasks, and make operational decisions 
under vague natural-language instructions. 
A future software system may optimize a user's calendar,
plan and rebook activities, or identify and execute 
financial investment decisions according to user 
preferences without receiving a fully rigorous 
specification of every intermediate step. 
Intelligence is therefore becoming an essential part of software systems, 
because it allows them to operate even when user 
requests are underspecified, ambiguous, or expressed 
in natural language rather than in precise specifications.

This trend of emergent software is exciting and promising, 
but it also changes the security 
problem fundamentally. Once software is expected to act under 
vague instructions and open-ended interaction, its working environments can no 
longer be assumed to be fully trusted or fully well-defined. 
The system must interpret arbitrary user inputs, external signals, 
and contextual information that may be incomplete, misleading, or 
adversarial. In addition, many emerging software systems now rely on 
large language models (LLMs) or related learning-based components 
whose internal behavior is partially opaque. 
Unlike conventional software logic that is explicitly written 
and can often be inspected or reasoned about directly, 
these components are not fully predictable in advance. 
When such systems operate in high-stakes settings, 
especially those involving financial assets, 
code execution, or sensitive data, failures can quickly 
translate into financial loss, unsafe actions, or information leakage.

This dissertation studies how to defend such high-stakes software systems 
in untrusted environments. Two families of systems are central in the thesis. 
The first is smart contracts in decentralized finance (DeFi), 
that is, blockchain-based financial services such as trading, 
lending, and asset exchange, which now represent a large and economically 
important class of deployed autonomous software~\cite{DeFillama}. 

Smart contracts already embody one form of software autonomy: once deployed, they execute financial logic directly on-chain and must accept transactions from arbitrary users, including strategic attackers. The second family is LLM-based coding systems, including coding assistants and coding agents that can search for information, generate code, and sometimes take actions on a machine. These systems represent a second form of software autonomy: they translate vague human goals into concrete code and system actions, while relying on opaque model behavior and weakly trusted online information such as tutorials, package pages, documentation, and search results.

The common challenge across these families is that harmful behavior may be induced by factors that are difficult to enumerate or validate in advance: adversarial interactions, unexpected execution patterns, opaque training data, or misleading external content. This shared setting motivates the structure of the thesis. The first two works study how to constrain high-stakes execution at runtime for smart contracts already exposed to public, adversarial environments. The first line of work is based on the observation that malicious transactions often exhibit abnormal behaviors compared with normal transactions, for example in asset movements, invocation patterns, or state updates. This motivates runtime invariants, namely rules that should continue to hold during benign executions. The second line of work is based on a complementary observation: many attacks exploit unexpected control flows across functions and across contracts, even when each local component may appear harmless in isolation. This motivates control-flow integrity defenses that validate whether a transaction follows legitimate interaction patterns.

The latter two works shift to LLM-centered coding systems. Layer III asks whether production models already emit harmful artifacts in ordinary developer workflows. Layer IV asks whether weakly trusted online information can poison coding agents before code is even deployed. Taken together, the dissertation develops and evaluates defenses across multiple stages of modern software operation. Some defenses directly constrain execution when software is already acting in a high-stakes environment. Others reveal or measure harmful behavior induced by untrusted information sources upstream. Under this view, smart contracts and coding agents are not two unrelated topics. They are two important instances of emerging software systems that increasingly act with autonomy in environments that cannot be fully trusted.
